\hypertarget{chapter-4-advanced-functions-and-callbacks}{%
\section{CHAPTER 4: ADVANCED FUNCTIONS AND
CALLBACKS}\label{chapter-4-advanced-functions-and-callbacks}}

\hypertarget{advanced-functions-callbacks}{%
\section{ADVANCED FUNCTIONS \&
CALLBACKS}\label{advanced-functions-callbacks}}

\hypertarget{advanced-functions}{%
\section{ADVANCED FUNCTIONS}\label{advanced-functions}}

\hypertarget{variadic-functions}{%
\subsection{2.1 Variadic Functions}\label{variadic-functions}}

\begin{lstlisting}[language={C++}]
#include <iostream>
#include <cstdarg>
using namespace std;

// Function with variable number of arguments
int sum(int count, ...) {
    va_list args;
    va_start(args, count);

    int total = 0;
    for (int i = 0; i < count; i++) {
        total += va_arg(args, int);
    }

    va_end(args);
    return total;
}

int main() {
    cout << sum(3, 10, 20, 30) << endl;     // 60
    cout << sum(5, 1, 2, 3, 4, 5) << endl;  // 15

    return 0;
}
\end{lstlisting}

\hypertarget{function-recursion}{%
\subsection{2.2 Function Recursion}\label{function-recursion}}

\begin{lstlisting}[language={C++}]
#include <iostream>
using namespace std;

// Factorial using recursion
int factorial(int n) {
    if (n <= 1) {
        return 1;  // Base case
    }
    return n * factorial(n - 1);  // Recursive case
}

// Fibonacci using recursion (inefficient)
int fibonacci(int n) {
    if (n <= 1) return n;
    return fibonacci(n - 1) + fibonacci(n - 2);
}

// Fibonacci with memoization (efficient)
int fib_memo(int n, int memo[]) {
    if (n <= 1) return n;
    if (memo[n] != -1) return memo[n];

    memo[n] = fib_memo(n - 1, memo) + fib_memo(n - 2, memo);
    return memo[n];
}

int main() {
    cout << factorial(5) << endl;  // 120
    cout << fibonacci(10) << endl; // 55

    int memo[11];
    for (int i = 0; i < 11; i++) memo[i] = -1;
    cout << fib_memo(10, memo) << endl;  // 55

    return 0;
}
\end{lstlisting}

\hypertarget{inline-functions}{%
\subsection{2.3 Inline Functions}\label{inline-functions}}

\begin{lstlisting}[language={C++}]
#include <iostream>
using namespace std;

// Inline function - compiler may expand code
inline int square(int x) {
    return x * x;
}

// Inline with condition
inline double max_value(double a, double b) {
    return (a > b) ? a : b;
}

int main() {
    cout << square(5) << endl;      // 25
    cout << max_value(3.5, 2.1) << endl;  // 3.5

    return 0;
}
\end{lstlisting}

\begin{center}\rule{0.5\linewidth}{0.5pt}\end{center}

\hypertarget{professional-notes-functional-depth}{%
\subsubsection{Professional Notes: Functional
Depth}\label{professional-notes-functional-depth}}

\hypertarget{recursion-and-tail-call-optimization-tco}{%
\paragraph{1. Recursion and Tail-Call Optimization
(TCO)}\label{recursion-and-tail-call-optimization-tco}}

Recursion is a technique where a function calls itself. * \textbf{Base
Case}: The condition where recursion stops. * \textbf{Stack Overflow}:
Deep recursion can exhaust the stack. * \textbf{Tail-Call Optimization}:
If the recursive call is the \emph{last} action in the function, some
compilers can transform it into a loop, saving stack space.

\begin{lstlisting}[language={C++}]
// Tail-recursive factorial
int factorial_tail(int n, int acc = 1) {
    if (n <= 1) return acc;
    return factorial_tail(n - 1, n * acc);
}
\end{lstlisting}

\hypertarget{callable-objects-functors}{%
\paragraph{2. Callable Objects
(Functors)}\label{callable-objects-functors}}

In C++, anything that can be invoked with \passthrough{\lstinline!()!}
is a callable. * \textbf{Function Pointers}:
\passthrough{\lstinline!void (*ptr)(int)!}. * \textbf{Functors}: Classes
that overload \passthrough{\lstinline!operator()!}. * \textbf{Lambdas
(C++11)}: Anonymous functions. *
\textbf{\passthrough{\lstinline!std::function!} (C++11)}: A polymorphic
wrapper for any callable.

\hypertarget{argument-dependent-lookup-adl---recap}{%
\paragraph{3. Argument Dependent Lookup (ADL) -
Recap}\label{argument-dependent-lookup-adl---recap}}

Functions are found in the namespaces of their arguments. This is why
you can call \passthrough{\lstinline!std::cout << obj!} without
qualifying the operator if it's defined in the same namespace as
\passthrough{\lstinline!obj!}.

\begin{center}\rule{0.5\linewidth}{0.5pt}\end{center}

\hypertarget{static-functions}{%
\subsection{2.4 Static Functions}\label{static-functions}}

\begin{lstlisting}[language={C++}]
#include <iostream>
using namespace std;

// File scope - only visible in this file
static void internal_function() {
    cout << "Internal function" << endl;
}

// Static with counter
int get_call_count() {
    static int count = 0;  // Persists between calls
    return ++count;
}

int main() {
    cout << get_call_count() << endl;  // 1
    cout << get_call_count() << endl;  // 2
    cout << get_call_count() << endl;  // 3

    internal_function();  // OK in same file

    return 0;
}
\end{lstlisting}

\begin{center}\rule{0.5\linewidth}{0.5pt}\end{center}

\hypertarget{function-pointers-callbacks}{%
\section{FUNCTION POINTERS \&
CALLBACKS}\label{function-pointers-callbacks}}

\hypertarget{function-pointers}{%
\subsection{3.1 Function Pointers}\label{function-pointers}}

\begin{lstlisting}[language={C++}]
#include <iostream>
using namespace std;

// Function pointer declaration: return_type (*name)(parameters)
int add(int a, int b) {
    return a + b;
}

int subtract(int a, int b) {
    return a - b;
}

int multiply(int a, int b) {
    return a * b;
}

int main() {
    // Declare function pointer
    int (*operation)(int, int);

    // Assign function to pointer
    operation = add;
    cout << operation(5, 3) << endl;  // 8

    operation = subtract;
    cout << operation(5, 3) << endl;  // 2

    operation = multiply;
    cout << operation(5, 3) << endl;  // 15

    return 0;
}
\end{lstlisting}

\hypertarget{function-pointer-arrays}{%
\subsection{3.2 Function Pointer Arrays}\label{function-pointer-arrays}}

\begin{lstlisting}[language={C++}]
#include <iostream>
using namespace std;

int add(int a, int b) { return a + b; }
int subtract(int a, int b) { return a - b; }
int multiply(int a, int b) { return a * b; }
int divide(int a, int b) { return a / b; }

int main() {
    // Array of function pointers
    int (*operations[4])(int, int) = {
        add, subtract, multiply, divide
    };

    int a = 20, b = 4;

    for (int i = 0; i < 4; i++) {
        cout << "Result: " << operations[i](a, b) << endl;
    }
    // Output: 24, 16, 80, 5

    return 0;
}
\end{lstlisting}

\hypertarget{callbacks}{%
\subsection{3.3 Callbacks}\label{callbacks}}

\begin{lstlisting}[language={C++}]
#include <iostream>
#include <vector>
using namespace std;

// Callback function type
typedef void (*Callback)(const string&);

class Button {
private:
    Callback on_click;

public:
    void set_click_handler(Callback callback) {
        on_click = callback;
    }

    void click() {
        if (on_click) {
            on_click("Button clicked!");
        }
    }
};

void handle_click(const string& message) {
    cout << message << endl;
}

int main() {
    Button button;
    button.set_click_handler(handle_click);
    button.click();  // Output: Button clicked!

    return 0;
}
\end{lstlisting}

\begin{center}\rule{0.5\linewidth}{0.5pt}\end{center}

\hypertarget{chapter-4-floating-point-bit-manipulation}{%
\section{Chapter 4: Floating Point \& Bit
Manipulation}\label{chapter-4-floating-point-bit-manipulation}}

Understanding how data is stored at the bit level and how floating-point
numbers approximate real values is essential for high-performance C++
engineering.

\hypertarget{floating-point-arithmetic}{%
\subsection{4.1 Floating Point
Arithmetic}\label{floating-point-arithmetic}}

Floating point numbers represent approximations of their assigned
values. This is due to the binary representation of base-10 decimals.

\hypertarget{the-imprecision-pitfall}{%
\subsubsection{1. The Imprecision
Pitfall}\label{the-imprecision-pitfall}}

A common mistake is assuming floating point equality will work as
expected:

\begin{lstlisting}[language={C++}]
double a = 0.1;
double b = 0.2;
double c = 0.3;
if (a + b == c) {
    std::cout << "Exact match!" << std::endl;
} else {
    std::cout << "Imprecise!" << std::endl; // Usually prints this
}
\end{lstlisting}

In IEEE 754 standard (used by most C++ compilers), 0.1 and 0.2 cannot be
represented exactly in binary, leading to a small rounding error when
added.

\hypertarget{comparison-strategy}{%
\subsubsection{2. Comparison Strategy}\label{comparison-strategy}}

Never use \passthrough{\lstinline!==!} with floats. Use an ``epsilon''
(a small tolerance):

\begin{lstlisting}[language={C++}]
#include <cmath>
#include <limits>

bool nearly_equal(double a, double b) {
    return std::abs(a - b) < std::numeric_limits<double>::epsilon();
}
\end{lstlisting}

\hypertarget{special-values-nan-and-infinity}{%
\subsubsection{3. Special Values: NaN and
Infinity}\label{special-values-nan-and-infinity}}

\begin{itemize}
\tightlist
\item
  \textbf{\passthrough{\lstinline!std::numeric\_limits<double>::quiet\_NaN()!}}:
  Not a Number (e.g., 0/0).
\item
  \textbf{\passthrough{\lstinline!std::numeric\_limits<double>::infinity()!}}:
  Infinity (e.g., 1/0).
\end{itemize}

\begin{center}\rule{0.5\linewidth}{0.5pt}\end{center}

\hypertarget{bitwise-mastery-low-level-optimization}{%
\subsection{4.2 Bitwise Mastery (Low-Level
Optimization)}\label{bitwise-mastery-low-level-optimization}}

Bitwise operators manipulate individual bits. Essential for embedded
systems, graphics, and cryptography.

\hypertarget{the-operators}{%
\subsubsection{The Operators}\label{the-operators}}

\begin{itemize}
\tightlist
\item
  \passthrough{\lstinline!\&!} (AND): Both bits must be 1.
\item
  \passthrough{\lstinline!|!} (OR): At least one bit must be 1.
\item
  \passthrough{\lstinline!\^!} (XOR): Bits must be different.
\item
  \passthrough{\lstinline!\~!} (NOT): Flip all bits.
\item
  \passthrough{\lstinline!<<!} (Left Shift): Multiply by 2\^{}N.
\item
  \passthrough{\lstinline!>>!} (Right Shift): Divide by 2\^{}N.
\end{itemize}

\hypertarget{god-tier-tricks}{%
\subsubsection{God-Tier Tricks}\label{god-tier-tricks}}

\begin{enumerate}
\def\labelenumi{\arabic{enumi}.}
\tightlist
\item
  \textbf{Check Odd/Even}: \passthrough{\lstinline!(x \& 1) == 0!}
  (Even). Faster than \passthrough{\lstinline!\% 2!}.
\item
  \textbf{Multiply by 2}: \passthrough{\lstinline!x << 1!}.
\item
  \textbf{Divide by 2}: \passthrough{\lstinline!x >> 1!}.
\item
  \textbf{Clear Last Set Bit}: \passthrough{\lstinline!x \& (x - 1)!}.
  Used to count set bits (Kernighan's Algorithm).
\item
  \textbf{Check Power of 2}:
  \passthrough{\lstinline!(x > 0) \&\& ((x \& (x - 1)) == 0)!}.
\item
  \textbf{Toggle Bit N}: \passthrough{\lstinline!x \^= (1 << N)!}.
\item
  \textbf{Set Bit N}: \passthrough{\lstinline!x |= (1 << N)!}.
\item
  \textbf{Clear Bit N}: \passthrough{\lstinline!x \&= \~(1 << N)!}.
\end{enumerate}

\begin{lstlisting}[language={C++}]
// Fast Power of 2 check
bool isPowerOf2(int x) {
    return x && !(x & (x - 1));
}
\end{lstlisting}

\begin{center}\rule{0.5\linewidth}{0.5pt}\end{center}

\hypertarget{bit-fields}{%
\subsection{4.3 Bit Fields}\label{bit-fields}}

Bit fields allow you to specify the number of bits each member of a
struct or class should occupy. This is crucial for matching hardware
protocols or saving space.

\begin{lstlisting}[language={C++}]
struct HardwareRegister {
    unsigned int enable : 1;  // 1 bit
    unsigned int mode   : 3;  // 3 bits
    unsigned int value  : 4;  // 4 bits
};
\end{lstlisting}

\textbf{Professional Note}: The exact layout of bit fields is
implementation-defined and depends on the platform's endianness and
alignment rules.

\begin{center}\rule{0.5\linewidth}{0.5pt}\end{center}

\hypertarget{professional-notes-bits-floats}{%
\subsubsection{Professional Notes: Bits \&
Floats}\label{professional-notes-bits-floats}}

\hypertarget{bit-manipulation-hacks}{%
\paragraph{1. Bit Manipulation Hacks}\label{bit-manipulation-hacks}}

\begin{itemize}
\tightlist
\item
  \textbf{Swapping without temp}:
  \passthrough{\lstinline!a \^= b; b \^= a; a \^= b;!} (Warning: Slower
  than a temp variable on modern CPUs due to pipeline stalls).
\item
  \textbf{Absolute value without branching}:
  \passthrough{\lstinline!(x + (x >> 31)) \^ (x >> 31)!} for 32-bit
  signed integers.
\end{itemize}

\hypertarget{floating-point-models}{%
\paragraph{2. Floating Point Models}\label{floating-point-models}}

\begin{itemize}
\tightlist
\item
  \textbf{\passthrough{\lstinline!fast-math!}}: A compiler flag (e.g.,
  \passthrough{\lstinline!-ffast-math!} in GCC) that allows the compiler
  to ignore some IEEE 754 rules for speed, potentially breaking
  precision.
\item
  \textbf{\passthrough{\lstinline!long double!}}: On many systems, this
  provides 80-bit or 128-bit precision for high-precision scientific
  calculations.
\end{itemize}
